\begin{textsample}{POS Dim 1 – rn – Score 40.00 – rn/rn\_em\_22.txt}  \label{ex:f1_pos_266}
6.2.
A ÁREA DE LINGUAGENS E SUAS \textbf{TECNOLOGIAS} A Área de Linguagens no Ensino Médio, no Estado do Rio Grande do Norte, compreende os componentes curriculares : Língua Portuguesa, Arte, Educação Física, Língua \textbf{Inglesa} e Língua \textbf{Espanhola}.
Estes componentes privilegiarão as múltiplas linguagens, posto que englobam \textbf{letramento} e \textbf{multiletramento} ; gêneros \textbf{discursivos} da cultura \textbf{digital} ; multimodalidade.
Não se trata, pois, de uma valorização da língua em si, mas do modo como a língua, como linguagem, se \textbf{materializa} no uso, nas diferentes práticas e campos de atuação social ; e \textbf{materializada} pelos gêneros \textbf{discursivos}.
Nesse sentido, a língua como linguagem verbal é parte dessas múltiplas linguagens.
Ainda, registre -se que a língua/linguagem são compreendidas como fenômenos sociais, históricos e, portanto, de caráter variável, heterogêneo e sensível aos contextos de uso, de modo que garanta os direitos linguísticos aos diferentes povos e grupos sociais presentes no Estado, com o compromisso de trabalhar pela manutenção e pela valorização de suas línguas e dialetos, marcas distintivas de seu patrimônio imaterial.
Os marcos legais que situam a Área de Linguagens e suas Tecnologias, enquanto oferta no ensino \textbf{médio}, são consubstanciados na Lei de Diretrizes e Bases da Educação Nacional ( no. 9394/96 ), em seu art. 35 -A, e no art. 11 da Resolução n. 3, de 21 de novembro de 2018, que atualiza as Diretrizes Curriculares Nacionais para o Ensino Médio ( doravante DCNEM ), compõem o conjunto de áreas que desenvolverão a Formação Geral Básica do estudante nessa etapa de ensino.
Ainda, as DCNEM, no Art. 11, § 4o e incisos I, IV, V e IX, apontam que a organização da Área deverá contemplar estudos e práticas de língua portuguesa, arte, educação física e língua \textbf{inglesa}.
Este último, informa que além da oferta de Língua \textbf{Inglesa} há a possibilidade de oferta de outras línguas estrangeiras.
No estado do Rio Grande do Norte, há o compromisso com a oferta da Língua \textbf{Espanhola} como a segunda língua estrangeira, sendo que esse componente apresenta -se inserido nos Itinerários Formativos, enquanto unidade curricular fixa.
Nesse caso, o deslocamento da carga horária do componente de Língua \textbf{Espanhola} para os Itinerários Formativos, se deu em virtude da legislação vigente ter estabelecido a carga horária máxima em 1.800h referente à Formação Geral Básica.
No entanto, enquanto orientação pedagógica, possui \textbf{matriz} específica integrada e articulada com os demais componentes na área de linguagens e suas \textbf{tecnologias}.
Desse modo, pautando -se nas concepções de Educação adotadas pelo Rio Grande do Norte e nas normativas apresentadas pela Base Nacional Comum Curricular - BNCC para a elaboração do Currículo como documento norteador das ações pedagógicas, a Área de Linguagens e suas Tecnologias buscará desenvolver nos estudantes \textbf{competências} e habilidades que promovam a articulação dos conhecimentos apreendidos em cada componente, de forma interdisciplinar, por exemplo, com integração e complementaridade entre os objetos de conhecimento de diferentes componentes da própria Área e transdisciplinar que, como uma estratégia, poderia tratar de um tema comum ( \textbf{transversal} ), ultrapassando os limites dos componentes para interagir com outros das demais Áreas do Conhecimento.
Ao mesmo tempo explorar as dimensões \textbf{socioemocionais} em situações de aprendizagem significativas e \textbf{relevantes} para a constituição do estudante, enquanto sujeito de direitos e deveres na interação em sociedade.
A partir do entendimento supramencionado, o objetivo da Área é proporcionar aos \textbf{alunos} do Ensino Médio ( EM ) a oportunidade de se apropriarem das diversas linguagens para que possam desenvolver uma visão \textbf{crítica}, ética e estética de seus usos oral, escrito, corporal, visual, sonoro e \textbf{digital}, manifestadas através de imagens, objetos artísticos visuais, gestos, música, teatro, movimentos corporais expressos pela dança e pelas atividades físicas, entre outras, bem como sua circulação nas Tecnologias \textbf{Digitais} de Informação e Comunicação ( TDIC ).
Destarte, tal postura \textbf{crítica} será incentivada \textbf{tanto} nos processos de pesquisa, na pré-iniciação científica e na busca de informação, como nas práticas de produção de sentido em seus diferentes componentes curriculares ( Arte, Educação Física, Língua Portuguesa, Língua \textbf{Inglesa} e Língua \textbf{Espanhola} ) na escola e, além dela, nos diversos campos da vida social, inclusive no processo de preparação para o mundo do trabalho.
A concepção de linguagem assumida neste documento é a enunciativa-discursiva, como posto em Bakthin ( 2010 ), que considera o discurso uma prática social e uma forma de interação, privilegiando a relação interpessoal, o contexto de produção dos textos, as diferentes situações de comunicação, os gêneros \textbf{discursivos}, a interpretação e a intenção de quem o produz, adequadas a cada campo da atividade humana.
Nessa perspectiva, o sujeito do discurso apresenta -se constituído na e pela linguagem, como detentor de subjetividade, perpassada por diversas \textbf{vozes} \textbf{discursivas}, sujeito \textbf{dialógico} por natureza.
Sendo assim, o Referencial Curricular do RN concebe os textos como pertencentes a gêneros \textbf{discursivos}, vinculados a diferentes contextos de produção e de \textbf{recepção}, que circulam em diversos campos sociais de atuação : Vida Pessoal ( VP ), Práticas de Estudo e Pesquisa ( PEP ), Jornalístico-Midiático ( JM ), Atuação na Vida Pública ( AVP ) e Campo Artístico ( CA ).
Dessa \textbf{maneira}, estará volvendo o olhar sobre as múltiplas culturas juvenis ou muitas \textbf{juventudes}, em constante (re)construção em face das novas exigências da \textbf{contemporaneidade} e o papel da escola na formação dos jovens, com \textbf{vistas} à construção de uma sociedade democrática e inclusiva.
Há também a preocupação de garantir que estes jovens permaneçam na escola, ou seja, que usufruam do direito público de todo cidadão brasileiro : completar a Educação Básica.
Nesse sentido, as DCN informam que os sujeitos do EM - adolescentes e jovens – não devem ser concebidos como um grupo homogêneo.
Por isso, este Referencial Curricular, em convergência, pretende respeitar as características dos jovens do Estado do Rio Grande do Norte ( RN ).
Nessa perspectiva, são consideradas as diversas produções que constituem as culturas juvenis manifestadas, por exemplo, em músicas populares como a MPB, regionais como o cordel, as canções de viola, a embolada, o forró e de \textbf{influência} internacional como o Hip-hop, o RAP, o slam, o funk, o Rock, dentre outras.
Ainda, nas danças, na cultura corporal, \textbf{autoria} de produções audiovisuais, rádios comunitárias, redes de mídia da internet, grafitagem, gírias e diferentes modos de organização da \textbf{juventude} norte-rio-grandense, de forma que lhes amplie o espaço destinado ao desenvolvimento de autonomia, protagonismo e \textbf{autoria} em situações reais de uso das diferentes práticas de linguagens.
De acordo com a BNCC do Ensino Médio, a Área de Linguagens e suas Tecnologias > [... ] tem a responsabilidade de propiciar oportunidades para a consolidação e a ampliação das habilidades de uso e de reflexão sobre as linguagens – artísticas, corporais e verbais ( oral ou \textbf{visual-motora}, como \textbf{Libras}, e escrita ) –, que são objeto de seus diferentes componentes ( Arte, Educação Física, Língua \textbf{Inglesa} e Língua Portuguesa ) ( BNCC-EM, p. 474 ).
Não obstante a BNCC não inclua Língua \textbf{Espanhola} como componente, o Referencial Curricular do Rio Grande do Norte o acolhe em respeito ao seu grande uso no Estado pela expansão do turismo e o desenvolvimento \textbf{econômico}, entendendo que outras línguas \textbf{modernas} podem ser incluídas, de acordo com as opções, necessidades e possibilidades de cada região.
Na direção de corresponder às \textbf{demandas} citadas, a \textbf{abordagem} pedagógica aqui \textbf{sugerida} \textbf{remete} a práticas de linguagens ( \textbf{compreensão}, interpretação e produção de textos \textbf{multissemióticos} em diferentes mídias e \textbf{análise} de suas formas de composição e estilo ), nos diversos campos/esferas de atuação humana, que, além de preparar os estudantes para lidarem com a informação cada vez mais acessível, nos mais diferentes contextos, tem como meta propor -lhes situações de aprendizagem que promovam o convívio com as diferenças e as diversidades, numa perspectiva de Educação em/para os Direitos Humanos, como também uma educação para a paz, visando à sustentabilidade.
Observamos que, dentre os componentes da área, apenas Língua Portuguesa apresenta carga horária prevista nos três anos do Ensino Médio, na Formação Geral Básica.
Para os demais, nessa parte comum, a carga horária é distribuída de \textbf{maneiras} diferentes ao longo do ensino \textbf{médio}.
Por essa \textbf{razão}, este Referencial propõe \textbf{matrizes} para cada componente com habilidades e objetos de conhecimento que deverão ser ministrados pelos professores em consonância com a estrutura curricular definida pela Secretaria Estadual de Educação do RN.
Assim, independentemente da distribuição de horas-aula, a Área de Linguagens possibilita aos estudantes oportunidades para vivenciarem práticas que os \textbf{aproximem} das reais \textbf{demandas} escolares, da vida pública, cultural e profissional e, ao mesmo tempo, propiciará crescimento pessoal, de modo a desenvolver \textbf{competências} nas suas formas de se organizarem, de construírem a própria identidade, de valorizarem sua participação na comunidade, articulando conhecimentos científicos com a vida prática.
Nesse sentido, a Área assume como \textbf{organizadores} da ação didática os eixos integradores e os campos de atuação social, conforme vemos a seguir :

% matched lemmas: abordagem, aluno, análise, aproximar, autoria, competência, compreensão, contemporaneidade, crítico, demanda, dialógico, digital, discursivo, econômico, espanhol, influência, inglês, juventude, letramento, libra, maneira, materializar, matriz, moderno, multiletramento, multissemiótico, médio, organizador, razão, recepção, relevante, remeter, socioemocional, sugerir, tanto, tecnologia, transversal, vista, visual-motor, voz
\end{textsample}
